% Additional main-paper theory used to fill the AAAI technical-content budget.

\subsection{Formal sufficient conditions}

The visible-spectrum definition is deliberately coordinate-free: it states what geometry the learned diagonal preconditioner realizes, without assuming how the feature coordinates were constructed.  The following results connect it to concrete feature maps.

\begin{proposition}[Band-limited mixing preserves visibility]
\label{prop:band-visibility}
Suppose the eigenvalues are partitioned into contiguous bands $B_1,\ldots,B_L$ with scales $\Lambda_\ell$ such that
\[
 \kappa^{-1}\Lambda_\ell\le\lambda_i(\Sigma)\le\kappa\Lambda_\ell,
 \qquad i\in B_\ell,
\]
for a scale-independent $\kappa$.  If the optimizer basis is obtained by an orthogonal rotation that is block diagonal over these bands, then the coordinatewise preconditioner satisfies
\[
 P\asymp_\kappa(\Sigma+\rho I)^{-1/2}.
\]
Consequently $\theta=1$, $\qeff=1/2$, and source mass is preserved exactly within every invariant band.
\end{proposition}

\begin{proof}[Proof sketch]
Within band $B_\ell$, every coordinate variance is a Rayleigh quotient of the corresponding covariance block and therefore lies in $[\kappa^{-1}\Lambda_\ell,\kappa\Lambda_\ell]$.  Hence the diagonal entries of $P$ are all within constants of $(\Lambda_\ell+\rho)^{-1/2}$.  Congruencing the block covariance by $P^{1/2}$ yields effective eigenvalues comparable to $\lambda_i(\lambda_i+\rho)^{-1/2}$.  Because both the covariance and preconditioner leave the band invariant, the quadratic source mass in that block is unchanged under the transformed-coordinate map.  Full details are in the supplement.
\end{proof}

\begin{proposition}[Flat diagonals hide spectral order]
\label{prop:flat-visibility}
Suppose $d_{\min}\le\Sigma_{jj}\le d_{\max}$ and $d_{\max}/d_{\min}\le\kappa$, where $\kappa$ is independent of scale.  Then
\[
 (d_{\max}+\rho)^{-1/2}\Sigma
 \pre P^{1/2}\Sigma P^{1/2}
 \pre(d_{\min}+\rho)^{-1/2}\Sigma.
\]
Thus the original and effective spectra have the same power-law exponent, so $\theta=0$ and $\qeff=0$.  The same exponent conclusion holds when $\kappa$ grows subpolynomially with model dimension.
\end{proposition}

The proposition explains why a change of basis can remove the exponent advantage without removing the useful covariance spectrum itself.  Diagonal adaptivity becomes scalar-comparable, while a full spectral preconditioner can still flatten the eigenvalues.

\begin{proposition}[A sufficient Gaussian-sketch condition]
\label{prop:gaussian-sketch}
Let the $M$ sketch rows be independent $s_j\sim\mathcal N(0,I_D/D)$ and let $d_j=s_j^\top Hs_j$.  Define
\[
 r_{\rm eff}=\frac{\tr(H)^2}{\tr(H^2)},
 \qquad
 r_{\rm st}=\frac{\tr(H)}{\|H\|_{\rm op}}.
\]
There is a universal $c>0$ such that, for $t\in(0,1)$,
\[
 \Pr\!\left(\max_{j\le M}
 \frac{|d_j-\tr(H)/D|}{\tr(H)/D}>t\right)
 \le2M\exp\{-c\min(t^2r_{\rm eff},t r_{\rm st})\}.
\]
Whenever the right-hand side tends to zero, the sketch-coordinate variances are uniformly flat and Proposition~\ref{prop:flat-visibility} applies.
\end{proposition}

This result is a sufficient condition, not a claim that every finite Gaussian sketch is perfectly flat.  In low effective-rank regimes, finite-dimensional fluctuations remain, and we directly measure the induced $\qeff$ rather than assigning a limiting value by assumption.  This distinction is important in the random-feature experiments below.

\begin{table}[t]
\centering
\small
\caption{Visible-spectrum regimes before the damping knee.}
\label{tab:regimes}
\begin{tabular}{lccc}
\toprule
Feature geometry & $\theta$ & $\qeff$ & Learned modes $K(n)$\\
\midrule
Aligned / band limited & $1$ & $1/2$ & $n^{2/a}$\\
Partial visibility & $(0,1)$ & $\theta/2$ & $n^{1/[a(1-\theta/2)]}$\\
Flat / global mixing & $0$ & $0$ & $n^{1/a}$\\
\bottomrule
\end{tabular}
\end{table}

\section{Further Consequences of Visible Spectrum}

\subsection{Time-limited and variance-limited phases}

The compute theorem has a useful piecewise interpretation.  Define the critical visibility
\[
 \theta_c=2\left(1-\frac ba\right).
\]
When $0<\theta_c<1$,
\begin{equation}
\label{eq:phase-transition}
 \beta(\theta)=
 \begin{cases}
 \dfrac{b-1}{a(1-\theta/2)+1}, & 0\le\theta\le\theta_c,\\[7pt]
 \dfrac{b-1}{b+1}, & \theta\ge\theta_c.
 \end{cases}
\end{equation}
Before $\theta_c$, optimization cannot expose modes as quickly as the variance-optimal rate and more visibility improves the compute exponent.  After $\theta_c$, bias and variance balance at $K\asymp N^{1/b}$; further flattening changes constants and the optimal scalar learning-rate schedule, but not the asymptotic compute exponent.  This explains why some of our compute sweeps improve monotonically over the full interval while others saturate early.

\subsection{Why tuning cannot replace spectral visibility}

The distinction between constants and exponents can be stated as an asymptotic separation result.  Let
\[
 \beta(q)=\frac{b-1}{\max\{a(1-q),b\}+1}
\]
be the compute exponent associated with an effective spectral exponent $q$.

\begin{proposition}[Hidden-spectrum oracle separation]
\label{prop:oracle-separation}
Suppose a coordinatewise method is scalar-comparable, so $q_{\rm coord}=0$, while a spectral oracle realizes $q_{\rm spec}=1/2$.  If $\beta(1/2)>\beta(0)$, then
\[
 \frac{R^\star_{\rm coord}(C)-\sigma^2}
 {R^\star_{\rm spec}(C)-\sigma^2}
 \asymp C^{\beta(1/2)-\beta(0)}\longrightarrow\infty.
\]
No scalar learning-rate retuning can remove this asymptotic gap.
\end{proposition}

\begin{proof}
The compute theorem gives $R_q^\star(C)-\sigma^2\asymp C^{-\beta(q)}$.  Taking the ratio proves the displayed result.  Replacing the effective horizon $n$ by $cn$ multiplies each spectral-filter term by a $c$-dependent constant but does not change its power in $n$ or $C$.
\end{proof}

This proposition motivates the spectral-oracle control in Figure~\ref{fig:experiments}.  In globally mixed coordinates, coordinatewise Adam can improve finite-time normalization while remaining in the $q=0$ class; the oracle reveals the performance left on the table when the same spectrum is available but hidden from a diagonal preconditioner.

\subsection{A measurable mechanism variable}

For an approximately power-law fitting window, let $\widehat a$ be the fitted exponent of $\Sigma$ and $\widehat\alpha$ the fitted exponent of $P^{1/2}\Sigma P^{1/2}$.  We use
\[
 \widehat q_{\rm eff}=1-\frac{\widehat\alpha}{\widehat a}.
\]
If the log spectra deviate from straight lines by at most $\eta$ over ranks $[L,U]$, ordinary least-squares perturbation gives
\[
 |\widehat q_{\rm eff}-q_{\rm eff}|
 \lesssim\frac{\eta}{\log(U/L)}
\]
when $a$ is bounded away from zero.  Thus $\qeff$ is not merely a visualization: it is a consistent diagnostic on widening power-law windows.  It also separates two questions that final loss alone cannot: whether an optimizer improves current performance, and whether it changes the asymptotic spectral class.
